\documentclass{article}

\usepackage[english]{babel}
\usepackage[ddmmyyyy]{datetime}
\usepackage[a4paper,top=2.54cm,bottom=2.54cm,left=2.54cm,right=2.54cm,marginparwidth=1.75cm]{geometry}

\title{Conference Abstract Proposal:\\Ductility and Fracture Toughness of Quasi-Distorted Metamaterials}
\author{Niccolo Forte}
\date{\today}

\begin{document}

\maketitle

\section*{}
Lattice structures are commonly implemented in a variety of engineering applications due to their ability to pertain significant aspects of their parent materials’ mechanical properties, while reducing overall structural weight. Such properties are of interest for various industries, including automotive, aerospace, and manufacturing, due to cost decreases incurred from reduced weight. Natural cellular materials, such as honeycombs, marine mussels, woods, trabecular bones, and sponges, may benefit from elements of disorderliness within their latticelike microstructures. Recent work has shown that introducing disorderliness into the geometry of such lattices, which was previously believed to adversely impact the mechanical properties of periodic truss structures (refs), can enhance their mechanical properties, leading to either ductile, damage tolerant, or catastrophic failure modes, independent of parent material properties (ref – AK2). Throughout this work we demonstrate how both ductility and fracture toughness can be enhanced through introduction of two forms of disorderliness, spatial coordinate perturbations and strut thickness variations, to perfectly periodic FCC, triangular, hexagonal, and kagome truss lattices.

To show parent material independence, numerical simulations using the Abaqus CAE explicit solver are conducted on random and patterned distributions of the two previously described forms of disorderliness in 2D models of the four lattice topologies for both a ductile, aluminium, and brittle, silicon carbide, material. Disorganised geometries resulting in the greatest enhancement of both the structure’s ductility and fracture toughness for each disorderliness and lattice type on both materials are additively manufactured using a ductile and brittle polymer to experimentally validate numerical results through plane-stress uniaxial tension and compact tension (CT) tests. The current work’s results will show that the mechanical properties of lattices are highly dependent on the distribution of disorderliness within the structure, which may result in various failure mods independent of the parent material’s mechanical properties. Some distributions may result in the sudden, catastrophic failure of the structure, whereas others may be tuned in such way to lead to a more desirable, progressive failure mode.

\end{document}